% !TeX program = pdflatex
% ---------------------------------------------------------------------------
% Equilibrium Risk Premia and Valuation under the Log-Normal SV Model
% Standalone working note, Zurich, August 2026.
%
% Provenance: this document is a corrected and extended version of the draft
% chapter ValuationUnderSV.tex. Blocks are marked as follows:
%   %% [CORRECTED vs chapter Eq (...)] - fixes an error in the draft
%   %% [NEW] - content added beyond the draft
% The tracked source snapshot does not contain verify_valuation_under_sv.py or
% verify_hawkes_equilibrium.py.  References to those maintainer working files are
% historical provenance, not claims that executable verification is present here.
% v2 (2026-08-23): added Section on the Hawkes-model equilibrium (companion to
% Liu-Packham-Sepp) and the discussion subsection on variance versus skew premia.
% v3 (2026-08-23): added Section on the regime-switching log-normal SV model
% (synthesis with the transition-tied jumps of the goal-based-allocation paper),
% derived and benchmarked in verify_regime_sv_equilibrium.py.
% v4 (2026-08-23): reconciled the regime notation and equity jump calibration
% with Sepp (2026), retained the induced cubic Q volatility drift, added the
% degree-four option transform, Monte Carlo validation, and smile sensitivities.
% v5 (2026-08-23): added the internally consistent log-linear equilibrium closure
% that preserves the quadratic Q drift, explained its relation to IJTAF, and
% compared its implied-volatility skews with the full-cubic closure.
% v6 (2026-08-23): clarified the exact and approximate senses in which the
% log-linear regime closure is consistent with the CRRA equilibrium framework.
% The economic-mechanism paragraph in Section 8 is an AI-drafted sketch and is
% marked for rewriting by the author, per the project drafting rule.
% ---------------------------------------------------------------------------
\documentclass[11pt]{article}
\usepackage[a4paper, margin=2.6cm]{geometry}
\usepackage{amsmath, amssymb, amsthm}
\usepackage{booktabs,float,graphicx}
\usepackage[round]{natbib}
\usepackage[colorlinks=true, linkcolor=blue, citecolor=blue]{hyperref}
% The default is relative to this tracked notes directory.  A parent book build
% may define \RegimeSvArtifactRoot before including this source.
\providecommand{\RegimeSvArtifactRoot}{../../../outputs/volatility_book/ch_lognormal_sv_risk_premia/canonical}
\graphicspath{{\RegimeSvArtifactRoot/figures/}}

\newtheorem{theorem}{Theorem}
\newtheorem{proposition}{Proposition}
\newtheorem{lemma}{Lemma}
\theoremstyle{remark}
\newtheorem{remark}{Remark}

\newcommand{\E}{\mathbb{E}}
\newcommand{\PP}{\mathbb{P}}
\newcommand{\QQ}{\mathbb{Q}}

\title{Equilibrium Risk Premia and Valuation\\ under the Log-Normal Stochastic Volatility Model\\[1ex] \large Working note}
\author{Artur Sepp}
\date{Zurich, August 2026}

\begin{document}
\maketitle

\begin{abstract}
\sloppy
\noindent
The log-normal stochastic volatility (SV) model with quadratic drift of \cite{SeppRakhmonov2024} specifies the market prices of equity and volatility risk as linear functions of the volatility. This specification is postulated for tractability rather than derived from economic primitives. We derive both risk premia in a production-economy equilibrium in which a representative agent with constant relative risk aversion (CRRA) holds the market portfolio. Under a local log-linear closure, both market prices of risk are proportional to volatility with coefficients $\bar\lambda_{0}=(1-\gamma)-\beta a^{(1)}$ and $\bar\lambda_{1}=-\varepsilon a^{(1)}$, where $a^{(1)}$ solves a scalar Riccati equation in closed form. Companion derivations for the bivariate Hawkes jump model and for a regime-switching extension show that the same equilibrium generates a common direct jump tilt $\gamma-1$. The regime-switching PDE for the value coefficients remains linear. Its log-linear closure gives four nonlinear ordinary differential equations and preserves the quadratic risk-neutral volatility drift, while its log-quadratic closure gives six equations and induces a cubic drift. A common degree-four transform yields semi-analytic option prices for both closures. We solve and simulate each internally consistent measure and find very small implied-volatility-skew differences for the equity illustration.
\end{abstract}

\section{Introduction}\label{sec:intro}

This note derives the equity and volatility risk premia of the log-normal SV model in general equilibrium. We follow the production-economy approach of \cite{CIR1985}: a representative CRRA agent allocates wealth to the stock index, a derivative security, and the money market account, and market clearing forces the agent to hold the index. The risk premia are outputs of the market-clearing conditions, not inputs. \cite{Heston1993} used the same argument to justify the volatility risk premium $\lambda V$ in the square-root model, and \cite{Lewis2000} develops it for general SV models. \cite{Bick1990} and \cite{HeLeland1993} characterize which price processes are consistent with such an equilibrium.

The note makes five points. First, the equilibrium delivers the market prices of risk in the factorized form $\lambda_{0}=(1-\gamma)\sigma-\beta(\sigma)\psi$ and $\lambda_{1}=-a(\sigma)\psi$, where $\psi$ is the logarithmic gradient of the value-function coefficient with respect to volatility. Second, the local log-linear solution makes both prices exactly proportional to $\sigma_{t}$ and provides closed-form coefficients. Third, this closure preserves the quadratic-drift LogSV class, whereas the more accurate log-quadratic equilibrium loading induces a cubic risk-neutral drift; the two statements refer to different approximation orders. Fourth, the Hawkes extension generates a common direct jump tilt $\xi^{+}=\xi^{-}=\gamma-1$. Fifth, the regime-switching LogSV model with the transition-tied jumps of \cite{Sepp2026gba} combines diffusive, regime-timing, and tail premia and remains semi-analytic after polynomial projection.

The note proceeds as follows. Section \ref{sec:model} states the model and the economy. Section \ref{sec:hjb} derives the Hamilton-Jacobi-Bellman (HJB) equation and the first-order conditions. Section \ref{sec:premia} imposes market clearing and derives the risk premia and the linear equation for the value-function coefficient. Section \ref{sec:pricing} states the equilibrium valuation equation. Section \ref{sec:lognormal} solves the log-normal case with the closed-form Riccati solution. Section \ref{sec:linearity} states the endogenous linear market prices of risk and the parameter mapping under the risk-neutral measure. Section \ref{sec:hawkes} extends the derivation to the bivariate Hawkes jump model and identifies its measure-change family with the equilibrium kernel. Section \ref{sec:regime} combines the two settings in a regime-switching log-normal SV model, derives both consistent equilibrium closures, and compares their implied-volatility skews. Section \ref{sec:discussion} discusses the horizon dependence of the premia, the economic mechanism, and the division of labor between diffusive and jump risk premia.

\section{Model and economy}\label{sec:model}

We specify the joint dynamics of the price process $S_{t}$ and the volatility process $\sigma_{t}$ under the statistical measure $\PP$ using the beta formulation of \cite{KarasinskiSepp2012}:
\begin{equation}\label{eq:sv}
\begin{split}
dS_{t} &= \mu_{t} S_{t} dt+ \sigma_{t}S_{t} dW^{(0)}_{t}, \quad S_0=S, \\
d\sigma_{t}&=b(\sigma_{t}) dt + \beta(\sigma_{t}) dW^{(0)}_{t} + a(\sigma_{t}) dW^{(1)}_{t} , \quad \sigma_{0}=\sigma,
\end{split}
\end{equation}
where $W^{(0)}_{t}$ and $W^{(1)}_{t}$ are independent Brownian motions, $\mu_{t}$ is the statistical drift, $b(\sigma)$ is the drift function of the volatility, $\beta(\sigma)$ is the volatility beta function, and $a(\sigma)$ is the residual volatility-of-volatility function. The beta function specifies the covariation between returns and volatility changes:
\begin{equation}\label{eq:covar}
\left\langle  \frac{dS_{t}}{S_{t}},\, d\sigma_{t} \right\rangle =  \sigma_{t} \beta(\sigma_{t})\, dt.
\end{equation}
The formulation \eqref{eq:sv} is equivalent to a classic two-factor SV model with correlated Brownian drivers. Given a volatility function $v(\sigma)$ and correlation $\rho$, we recover \eqref{eq:sv} by setting $\beta(\sigma)=\rho v(\sigma)$ and $a(\sigma)=\sqrt{1-\rho^{2}}\,v(\sigma)$.

%% [NEW] problem statement was missing in the draft
The economy consists of three assets: the stock index $S_{t}$ in unit net supply, the money market account $B_{t}$ with $dB_{t}=rB_{t}dt$ in zero net supply, and a derivative security in zero net supply with value function $U(t, S_{t}, \sigma_{t})$ and maturity $T_{U}$. A representative agent allocates the fractions $x_{t}$, $y_{t}$ and $z_{t}=1-x_{t}-y_{t}$ of wealth $\Pi_{t}$ to the index, the derivative, and the account, respectively. The agent maximizes the expected utility of terminal wealth at horizon $T$, $T \geq T_{U}$, with CRRA preferences:
\begin{equation}\label{eq:objective}
J(t, \Pi, \sigma) = \sup_{ \{x_{s}, y_{s}\}_{s \in [t,T]} } \E\left[ \left. \frac{\Pi_{T}^{\gamma}}{\gamma} \right| \Pi_{t}=\Pi, \sigma_{t}=\sigma \right], \quad \gamma<1, \ \gamma \neq 0,
\end{equation}
so that the relative risk aversion equals $1-\gamma$. There is no intermediate consumption. In equilibrium the agent holds the net supply of each asset, which requires $x^{*}=1$, $y^{*}=0$ and $z^{*}=0$, and the drifts $\mu_{t}$ and $\mu^{U}_{t}$ adjust to make these allocations optimal.

Table \ref{tab:notation} collects the notation used throughout the note.
\begin{table}[h]
\centering
\begin{tabular}{ll}
\hline
$b(\sigma)$, $\beta(\sigma)$, $a(\sigma)$ & drift, beta, and residual volatility functions of $\sigma_{t}$ \\
$\gamma$ & CRRA utility power, risk aversion $1-\gamma$ \\
$J(t,\Pi,\sigma)$, $g(t,\sigma)$ & value function and its volatility coefficient, Eq \eqref{eq:ansatz} \\
$\psi(t,\sigma)=g_{\sigma}/g$ & risk-premium loading \\
$\lambda_{0}$, $\lambda_{1}$ & market prices of risk of $W^{(0)}$ and $W^{(1)}$ \\
$\varphi = \beta\lambda_{0} + a\lambda_{1}$ & drift adjustment of $\sigma_{t}$ under $\QQ$ \\
$\kappa_{1},\kappa_{2},\theta,\beta,\varepsilon$ & log-normal model parameters, Eq \eqref{eq:lnsv} \\
$\vartheta^{2}=\beta^{2}+\varepsilon^{2}$, $\bar\kappa = \kappa_{1}+\kappa_{2}\theta$ & total variance of volatility, linearized mean reversion \\
$\tau=T-t$, $v = \sigma-\theta$ & time to horizon, mean-adjusted volatility \\
\hline
\end{tabular}
\caption{Notation.}
\label{tab:notation}
\end{table}

\section{Wealth dynamics and the HJB equation}\label{sec:hjb}

Applying It\^{o}'s lemma to $U(t,S_{t},\sigma_{t})$ under the dynamics \eqref{eq:sv}, we obtain
\begin{equation}\label{eq:dU}
dU_{t} = \mu^{U}_{t} U_{t}\, dt+ \left[\sigma_{t}S_{t}\partial_{S}U + \beta(\sigma_{t})\partial_{\sigma}U\right] dW^{(0)}_{t}+ a(\sigma_{t})\, \partial_{\sigma}U\, dW^{(1)}_{t},
\end{equation}
where the drift rate $\mu^{U}_{t}$ is given by
%% [CORRECTED vs chapter Eq (vru2)]: the gamma term carries S^2
\begin{equation}\label{eq:muU}
\mu^{U} U = \partial_{t} U + \mu S \partial_{S} U + \frac{1}{2}\sigma^{2} S^{2} \partial^{2}_{SS} U  +  b(\sigma) \partial_{\sigma} U + \frac{1}{2}\left[ \beta^{2}(\sigma) + a^{2}(\sigma) \right] \partial^{2}_{\sigma\sigma} U + \sigma\beta(\sigma)S\,\partial^{2}_{S\sigma} U.
\end{equation}

Since $y_{t}$ is the fraction of wealth in the derivative, the number of derivative units is $y_{t}\Pi_{t}/U_{t}$, and the per-dollar exposures of the derivative to the two Brownian drivers are
%% [CORRECTED vs chapter Eq (log3)]: exposures are normalized by U
\begin{equation}\label{eq:deltas}
\Delta_{0} = \frac{\sigma S\,\partial_{S}U + \beta(\sigma)\,\partial_{\sigma}U}{U}, \qquad
\Delta_{1} = \frac{a(\sigma)\,\partial_{\sigma}U}{U}.
\end{equation}
The self-financing wealth process then follows
\begin{equation}\label{eq:wealth}
d\Pi_{t} = \left[(\mu_{t}-r)x_{t} + (\mu^{U}_{t} - r) y_{t} + r\right] \Pi_{t}\, dt
+ \left[ \sigma_{t}x_{t}  + y_{t} \Delta_{0} \right] \Pi_{t}\, dW^{(0)}_{t} + y_{t} \Delta_{1} \Pi_{t}\, dW^{(1)}_{t}.
\end{equation}

The value function \eqref{eq:objective} satisfies the HJB equation
\begin{equation}\label{eq:hjb}
\max_{x, y} \left\{ \partial_{t} J + \mathcal{L}^{(x,y)}(J)\right\} = 0, \quad J(T,\Pi,\sigma)=\frac{\Pi^{\gamma}}{\gamma},
\end{equation}
with the controlled generator
\begin{equation}\label{eq:generator}
\begin{split}
\mathcal{L}^{(x,y)}(J) & =  \left[(\mu-r)x + (\mu^{U} - r)y + r\right]\Pi J_{\Pi} + b(\sigma)J_{\sigma} +\frac{1}{2}\left[\beta^{2}(\sigma) + a^{2}(\sigma)\right]J_{\sigma\sigma}\\
& + \frac{1}{2}\left[ \left(\sigma x + y\Delta_{0}\right)^{2} + y^{2}\Delta_{1}^{2}\right]\Pi^{2} J_{\Pi\Pi}
+ \left[ \beta(\sigma) \left(\sigma x + y\Delta_{0}\right)  + a(\sigma)\, y \Delta_{1} \right] \Pi J_{\Pi\sigma}.
\end{split}
\end{equation}
Differentiating \eqref{eq:generator} in $x$ and $y$ and evaluating at the market-clearing allocations $x^{*}=1$ and $y^{*}=0$ yields the two first-order conditions
\begin{equation}\label{eq:focs}
\begin{split}
& (\mu-r) J_{\Pi} + \sigma^{2}\Pi J_{\Pi\Pi}+ \beta(\sigma)\, \sigma  J_{\Pi\sigma} = 0,\\
& (\mu^{U} - r) J_{\Pi} + \sigma \Delta_{0} \Pi J_{\Pi\Pi}+  \left[ \beta(\sigma) \Delta_{0} + a(\sigma)\Delta_{1} \right] J_{\Pi\sigma} = 0.
\end{split}
\end{equation}

\section{Equilibrium risk premia}\label{sec:premia}

We introduce the risk-aversion and hedging coefficients of the value function,
\begin{equation}\label{eq:AB}
A = - \frac{\Pi J_{\Pi \Pi}}{J_{\Pi}}, \qquad B= -\frac{J_{\Pi\sigma}}{J_{\Pi}}.
\end{equation}
The first condition in \eqref{eq:focs} determines the equity risk premium:
\begin{equation}\label{eq:equity_premium}
\mu-r = \sigma^{2}A + \beta(\sigma)\, \sigma B.
\end{equation}
Substituting \eqref{eq:equity_premium} and \eqref{eq:deltas} into the second condition in \eqref{eq:focs} and multiplying by $U$, we obtain the drift restriction on the derivative:
%% [CORRECTED vs chapter Eqs (log12), (log14)]: the left-hand side carries U
\begin{equation}\label{eq:U_premium}
U (\mu^{U} - r)  = (\mu-r)\,S\partial_{S}U + \varphi\, \partial_{\sigma}U,
\qquad
\varphi = \sigma\beta(\sigma) A + \left[ \beta^{2}(\sigma)  + a^{2}(\sigma) \right]B.
\end{equation}

Substituting the equilibrium premium \eqref{eq:equity_premium} back into \eqref{eq:generator} at $x^{*}=1$, $y^{*}=0$ reduces the HJB equation to
\begin{equation}\label{eq:Jpde}
\partial_{t} J + r\Pi J_{\Pi} - \frac{1}{2}\sigma^{2}\Pi^{2} J_{\Pi\Pi} +  b(\sigma)J_{\sigma} +\frac{1}{2}\left[\beta^{2}(\sigma) + a^{2}(\sigma)\right]J_{\sigma\sigma} = 0.
\end{equation}
For CRRA preferences we apply the multiplicative ansatz
\begin{equation}\label{eq:ansatz}
J(t,\Pi,\sigma) = g(t, \sigma)\, \frac{\Pi^{\gamma}}{\gamma},
\end{equation}
where we call $g(t,\sigma)$ the risk-premium coefficient. The ansatz gives $A = 1-\gamma$ and $B = -\psi(t,\sigma)$ with the risk-premium loading
\begin{equation}\label{eq:psi}
\psi (t, \sigma) = \frac{g_{\sigma}(t, \sigma) }{g(t, \sigma) },
\end{equation}
and \eqref{eq:Jpde} becomes a linear equation for $g$:
\begin{equation}\label{eq:gpde}
g_{t} +  b(\sigma)g_{\sigma} +\frac{1}{2}\left[\beta^{2}(\sigma) + a^{2}(\sigma)\right]g_{\sigma\sigma} +  \left[\gamma r  - \frac{1}{2} \gamma (\gamma-1) \sigma^{2} \right]g  = 0, \quad g(T, \sigma) = 1.
\end{equation}

\begin{remark}[No Merton distortion]\label{rem:linearity}
%% [NEW]
In the Merton problem with an exogenous drift $\mu$, optimizing over $x$ produces the nonlinear term $\frac{\gamma}{2(1-\gamma)}\left(\frac{\mu-r}{\sigma}\right)^{2}g$ in the equation for $g$, see \cite{Merton1971}. Here market clearing pins $x^{*}=1$ before the premium is known, and the premium \eqref{eq:equity_premium} adjusts to support this allocation. As a result, Eq \eqref{eq:gpde} is linear and contains no risk-premia inputs: the premia are outputs computed from its solution.
\end{remark}

%% [NEW] probabilistic representation and existence
By the Feynman-Kac theorem, the solution of \eqref{eq:gpde} admits the representation
\begin{equation}\label{eq:FK}
g(t,\sigma) = e^{\gamma r (T-t)}\, \E\left[ \left. \exp\left( \frac{1}{2}\gamma(1-\gamma) \int_{t}^{T} \sigma^{2}_{s}\, ds \right) \right| \sigma_{t}=\sigma \right],
\end{equation}
where the expectation is taken under $\PP$. The representation identifies $g$ with the moment generating function of the quadratic variation evaluated at the transform variable $\Psi = -\frac{1}{2}\gamma(1-\gamma)$, so the analytics developed for the quadratic variation in \cite{SeppRakhmonov2024} apply to $g$ directly.

\begin{proposition}[Existence]\label{prop:existence}
%% [NEW]
Let the volatility follow the log-normal dynamics \eqref{eq:lnsv} below. For $\gamma < 0$, the expectation \eqref{eq:FK} is finite for all $T$. For $0<\gamma<1$, it is finite for all $T$ provided that
\begin{equation}\label{eq:existence}
\kappa_{2} > \vartheta \sqrt{ \gamma(1-\gamma)}.
\end{equation}
\end{proposition}
\begin{proof}
For $\gamma<0$ the integrand exponent is negative, so the expectation is bounded by one. For $0<\gamma<1$ the claim follows from the supersolution argument in the proof of the existence theorem for the quadratic-variation moment generating function in \cite{SeppRakhmonov2024}, applied under $\PP$ with $-\Psi_{R} = \frac{1}{2}\gamma(1-\gamma)$.
\end{proof}

The market prices of risk of the two Brownian drivers follow from \eqref{eq:equity_premium} and from decomposing $\varphi$ in \eqref{eq:U_premium}:
%% [NEW] per-Brownian decomposition replaces the scalar lambda^sigma of the draft
\begin{equation}\label{eq:mprs}
\lambda_{0} = \frac{\mu-r}{\sigma} = (1 - \gamma)\sigma  - \beta(\sigma)\, \psi (t, \sigma),
\qquad
\lambda_{1} = -a(\sigma)\, \psi (t, \sigma),
\end{equation}
so that $\varphi = \beta(\sigma)\lambda_{0} + a(\sigma)\lambda_{1}$. The change of measure to the risk-neutral measure $\QQ$ is
\begin{equation}\label{eq:girsanov}
\left.\frac{d\QQ}{d\PP}\right|_{\mathcal{F}_t}=\mathcal{E}\left(-\int_0^t \lambda_0(s)\,dW^{(0)}_s-\int_0^t\lambda_1(s)\,dW^{(1)}_s\right),
\end{equation}
where $\mathcal{E}(\cdot)$ denotes the stochastic exponential. The residual price of risk $\lambda_{1}$ is proportional to the loading $\psi$ alone: the agent prices the orthogonal volatility shocks only through the sensitivity of the value function to the volatility state. For log utility, $\gamma \to 0$, the coefficient $g$ becomes independent of $\sigma$, so $\psi=0$, $\lambda_{0} = \sigma$ and $\lambda_{1} = 0$.

\section{Equilibrium valuation equation}\label{sec:pricing}

Equating the drift restriction \eqref{eq:U_premium} with the It\^{o} drift \eqref{eq:muU} yields the valuation equation for any derivative in zero net supply:
%% [CORRECTED vs chapter Eq (log24)]: the cross term carries S
\begin{equation}\label{eq:pricing_pde}
\begin{split}
&\partial_{t} U  + r S\partial_{S}U + \frac{1}{2} \sigma^{2} S^{2}\partial^{2}_{SS}U+ \left[b(\sigma) -\varphi (t, \sigma) \right]\partial_{\sigma}U \\
&\quad + \frac{1}{2}\left[\beta^{2}(\sigma) + a^{2}(\sigma)\right]\partial^{2}_{\sigma\sigma}U + \sigma\beta(\sigma) S\, \partial^{2}_{S\sigma}U - rU = 0.
\end{split}
\end{equation}
Equation \eqref{eq:pricing_pde} is the Feynman-Kac equation for the $\QQ$-dynamics
\begin{equation}\label{eq:Qdynamics}
\begin{split}
dS_{t} &= r S_{t}\, dt+ \sigma_{t}S_{t}\, d\hat{W}^{(0)}_{t}, \\
d\sigma_{t}&=\left[ b(\sigma_{t}) - \varphi(t, \sigma_{t}) \right] dt + \beta(\sigma_{t})\, d\hat{W}^{(0)}_{t} + a(\sigma_{t})\, d\hat{W}^{(1)}_{t},
\end{split}
\end{equation}
where $\hat{W}^{(0)}_{t}$ and $\hat{W}^{(1)}_{t}$ are independent Brownian motions under $\QQ$ defined by the change of measure \eqref{eq:girsanov}. The computation of the premia therefore reduces to three steps: solve the linear equation \eqref{eq:gpde} for $g$, compute the loading $\psi = g_{\sigma}/g$, and evaluate $\lambda_{0}$, $\lambda_{1}$ and $\varphi$ from \eqref{eq:mprs}.

\section{Solution for the log-normal volatility model}\label{sec:lognormal}

We specialize to the log-normal volatility dynamics with quadratic drift:
\begin{equation}\label{eq:lnsv}
d\sigma_t=(\kappa_1+\kappa_2\sigma_t)(\theta-\sigma_t)\, dt+\beta\sigma_t\, dW^{(0)}_t+\varepsilon\sigma_t\, dW^{(1)}_t,
\end{equation}
so that $b(\sigma) = (\kappa_1+\kappa_2\sigma)(\theta-\sigma)$, $\beta(\sigma)=\beta\sigma$ and $a(\sigma)=\varepsilon\sigma$, with the total variance of the volatility denoted by $\vartheta^{2}=\beta^{2}+\varepsilon^{2}$. To remove the interest-rate term, we factor $g = e^{\gamma r (T-t)}\, \hat{g}$, and Eq \eqref{eq:gpde} becomes
%% [CORRECTED vs chapter Eq (lsv2)]: the gamma-r term is factored out explicitly
\begin{equation}\label{eq:gpde_ln}
\hat{g}_{t} +   (\kappa_1+\kappa_2\sigma)(\theta-\sigma)\, \hat{g}_{\sigma} +\frac{1}{2} \vartheta^{2} \sigma^{2}\, \hat{g}_{\sigma\sigma} - \frac{1}{2} \gamma (\gamma-1) \sigma^{2}\, \hat{g}  = 0, \quad \hat{g}(T, \sigma) = 1.
\end{equation}

We change variables to the mean-adjusted volatility $v=\sigma - \theta$ and time to horizon $\tau=T-t$. Expanding the drift around $\theta$ gives $(\kappa_1+\kappa_2\sigma)(\theta-\sigma) = -\bar\kappa v - \kappa_2 v^{2}$ with the linearized mean-reversion speed $\bar\kappa = \kappa_{1}+\kappa_{2}\theta$, and Eq \eqref{eq:gpde_ln} becomes
%% [CORRECTED vs chapter Eq (lsv3)]: the linear drift coefficient is -(kappa_1+kappa_2*theta),
%% not (kappa_1-kappa_2*theta); the latter is the coefficient of the expansion around sigma=0
\begin{equation}\label{eq:gpde_v}
- \hat{g}_{\tau} -   \left[\bar\kappa\, v + \kappa_2 v^{2}\right]\hat{g}_{v} +\frac{1}{2} \vartheta^{2} (v+\theta)^{2}\, \hat{g}_{vv} - \frac{1}{2} \gamma (\gamma-1) (v+\theta)^{2} \hat{g}  = 0, \quad \hat{g}(0, v) = 1.
\end{equation}
The coefficient $\bar\kappa$ agrees with the linearized mean reversion $\kappa^{(1)} = \kappa_{1}+\kappa_{2}\theta$ in the operator representation of \cite{SeppRakhmonov2024}.

We first seek a local log-linear closure
\begin{equation}\label{eq:affine}
\hat{g}(\tau, v) = \exp \left\{ a^{(0)}(\tau) + a^{(1)}(\tau)\, v \right\} + O(v^{2}), \quad a^{(0)}(0)=a^{(1)}(0)=0,
%% [CORRECTED vs chapter Eq (lsv4)]: initial conditions are a0(0)=0 and a1(0)=0
\end{equation}
where the $O(v^{2})$ terms collect the discarded powers. This degree-one log-polynomial is a useful closed-form leading approximation; it is not the ``FIRST'' affine expansion terminology of \cite{SeppRakhmonov2024}, whose exponent has degree two. Collecting the terms of orders $v^{0}$ and $v^{1}$ yields the system
\begin{equation}\label{eq:odes}
\begin{split}
& a^{(0)}_{\tau}  = \frac{1}{2} \vartheta^{2} \theta^{2} \left(a^{(1)}\right)^{2} - \frac{1}{2} \theta^{2}  \gamma (\gamma-1),\\
&  a^{(1)}_{\tau}  =  \vartheta^{2} \theta \left(a^{(1)}\right)^{2} -\bar\kappa\,  a^{(1)}  - \theta \gamma (\gamma-1).
\end{split}
\end{equation}
The equation for $a^{(1)}$ is a scalar Riccati equation with constant coefficients, and it admits a closed-form solution.

\begin{lemma}[Closed-form solution]\label{lem:riccati}
%% [NEW]: the draft solved only for the steady state
Define the discriminant and the roots
\begin{equation}\label{eq:roots}
d= \sqrt{ \bar\kappa^{2}+4\vartheta^{2} \theta^{2} \gamma (\gamma-1) }, \qquad
r_{\pm} = \frac{\bar\kappa \pm d}{2\vartheta^{2} \theta},
\end{equation}
and assume $d^{2}>0$. Then the solution of \eqref{eq:odes} is
\begin{equation}\label{eq:riccati_solution}
a^{(1)}(\tau) = \frac{ r_{+} r_{-} \left(1 - e^{-d \tau}\right)}{r_{+} - r_{-}\, e^{-d \tau}},
\qquad
a^{(0)}(\tau) = \frac{\theta}{2}\, a^{(1)}(\tau) + \frac{\theta \bar\kappa}{2} \left[ r_{-}\tau + \frac{1}{\vartheta^{2}\theta} \ln \frac{r_{+}-r_{-}}{r_{+}-r_{-}e^{-d\tau}} \right].
\end{equation}
As $\tau \to \infty$, the loading converges to the steady state $a^{(1)}_{\infty} = r_{-}$.
\end{lemma}
\begin{proof}
Direct substitution verifies that $a^{(1)}$ in \eqref{eq:riccati_solution} solves the Riccati equation with $a^{(1)}(0)=0$. For $a^{(0)}$, the identity $a^{(0)}_{\tau} = \frac{\theta}{2}\left[ a^{(1)}_{\tau} + \bar\kappa a^{(1)} \right]$ follows from \eqref{eq:odes}, and integrating $a^{(1)}$ in closed form gives the stated expression.
\end{proof}

\begin{remark}[Branch selection and validity]\label{rem:branch}
%% [CORRECTED vs chapter Eqs (lsv8), (lsv9)] and [NEW] justification
The steady state $r_{-}$ is the attracting fixed point of the Riccati flow started at $a^{(1)}(0)=0$, while $r_{+}$ is repelling. The root $r_{-}$ also vanishes in the two limits without volatility risk compensation, $\gamma \to 0$ and $\gamma \to 1$, while $r_{+}$ does not, which confirms the selection on economic grounds. For $0<\gamma<1$ the discriminant is positive if $\bar\kappa^{2} \geq 4\vartheta^{2}\theta^{2}\gamma(1-\gamma)$, for which $\bar\kappa \geq \vartheta\theta$ is sufficient since $\gamma(1-\gamma) \leq 1/4$. For $\gamma<0$ the discriminant is always positive.
\end{remark}

\section{Endogenous linear market prices of risk}\label{sec:linearity}

%% [NEW] this section is the main addition relative to the draft
Under the log-linear closure, the loading $\psi = \hat{g}_{v}/\hat{g} = a^{(1)}(\tau)$ does not depend on the volatility state. Substituting $\beta(\sigma)=\beta\sigma$ and $a(\sigma)=\varepsilon\sigma$ into \eqref{eq:mprs}, both market prices of risk become exactly proportional to the volatility:
\begin{equation}\label{eq:linear_mprs}
\lambda_{0}(t) = \bar\lambda_{0}(\tau)\, \sigma_{t}, \quad \bar\lambda_{0}(\tau) = (1-\gamma) - \beta\, a^{(1)}(\tau),
\qquad
\lambda_{1}(t) = \bar\lambda_{1}(\tau)\, \sigma_{t}, \quad \bar\lambda_{1}(\tau) = -\varepsilon\, a^{(1)}(\tau).
\end{equation}
The linear specification $\lambda_{i} = \bar\lambda_{i}\sigma$, which \cite{SeppRakhmonov2024} adopt with reference to \cite{BakshiKapadia2003}, is therefore not an auxiliary assumption in this economy: it is the form the equilibrium delivers. The equity risk premium equals $\mu - r = \bar\lambda_{0}\sigma^{2}$, which is proportional to the instantaneous variance.

The drift adjustment becomes $\varphi = \left[ (1-\gamma)\beta - \vartheta^{2} a^{(1)} \right]\sigma^{2}$, so the $\QQ$-drift of the volatility retains the quadratic form:
\begin{equation}\label{eq:Qdrift}
b(\sigma) - \varphi = \kappa_{1}\theta - \left(\kappa_{1}-\kappa_{2}\theta\right)\sigma - \hat\kappa_{2}\, \sigma^{2},
\qquad
\hat\kappa_{2} = \kappa_{2} + (1-\gamma)\beta - \vartheta^{2} a^{(1)},
\end{equation}
which agrees with $\hat\kappa_{2} = \kappa_{2}+\beta\bar\lambda_{0}+\varepsilon\bar\lambda_{1}$ under \eqref{eq:linear_mprs}. The remaining risk-neutral parameters follow from the mapping in \cite{SeppRakhmonov2024}:
\begin{equation}\label{eq:Qparams}
\hat \theta =  \frac{ - (\kappa_1-\kappa_2\theta) + \sqrt{(\kappa_1-\kappa_2\theta)^{2}+4\theta\kappa_{1}\hat\kappa_2}}{2\hat\kappa_2}, \qquad
\hat \kappa_1 = \frac{ (\kappa_1-\kappa_2\theta) +  \sqrt{(\kappa_1-\kappa_2\theta)^{2}+4\theta\kappa_{1}\hat\kappa_2}}{2},
\end{equation}
so that the $\QQ$-dynamics are $d\sigma_{t} = (\hat\kappa_{1}+\hat\kappa_{2}\sigma_{t})(\hat\theta - \sigma_{t})dt + \beta\sigma_{t}d\hat{W}^{(0)}_{t}+\varepsilon\sigma_{t}d\hat{W}^{(1)}_{t}$. The equilibrium thereby preserves the invariance of the model class under the change of measure, with all $\QQ$-parameters stated in terms of the $\PP$-parameters and the risk aversion $1-\gamma$.\footnote{A reference implementation of the premia $\bar\lambda_{0}$, $\bar\lambda_{1}$ and the parameter mapping \eqref{eq:Qdrift}-\eqref{eq:Qparams} is planned for the \texttt{stochvolmodels} package, which implements the log-normal SV model of \cite{SeppRakhmonov2024}.}

\section{Equilibrium jump risk premia in the Hawkes model}\label{sec:hawkes}

%% [NEW] companion section: makes the correspondence with the jump-premia paper precise.
%% Historical working notes named verify_hawkes_equilibrium.py, but that script is
%% not part of the tracked source snapshot.  This section therefore records the
%% derivation without claiming an executable Hawkes verification artifact here.
We extend the equilibrium derivation to the bivariate Hawkes jump model of \cite{LiuPackhamSepp2026}, where the jump risk premia are specified through a parametric change of measure with two external tilt parameters $\chi^{+}$ and $\chi^{-}$ calibrated to option prices. We show that the CRRA equilibrium generates exactly this family and restricts the direct jump tilts to a common value set by the risk aversion.

\subsection{Model and economy}

The price process under $\PP$ carries a diffusive component and two self- and cross-exciting jump components:
\begin{equation}\label{eq:hawkes_dS}
\begin{split}
\frac{dS_{t}}{S_{t-}} &= \mu_{t}\, dt + \sigma\, dW_{t} + \left[ (e^{j^{+}}-1)\, dN^{(1)}_{t} - \lambda^{+}_{t} m^{+} dt \right] + \left[ (e^{j^{-}}-1)\, dN^{(2)}_{t} - \lambda^{-}_{t} m^{-} dt \right],\\
d\lambda^{+}_{t} &= \kappa^{+}(\theta^{+}-\lambda^{+}_{t})\, dt + \beta_{11}\, j^{+} dN^{(1)}_{t} + \beta_{12}\, j^{-} dN^{(2)}_{t},\\
d\lambda^{-}_{t} &= \kappa^{-}(\theta^{-}-\lambda^{-}_{t})\, dt + \beta_{21}\, j^{+} dN^{(1)}_{t} + \beta_{22}\, j^{-} dN^{(2)}_{t},
\end{split}
\end{equation}
where $N^{(1)}_{t}$ and $N^{(2)}_{t}$ count positive and negative jumps with intensities $\lambda^{+}_{t}$ and $\lambda^{-}_{t}$, the jump sizes $j^{\pm}$ have densities $\varpi^{\pm}$, and $W_{t}$ is an independent Brownian motion. We write the moment functions of the jump sizes as
\begin{equation}\label{eq:hawkes_ell}
\ell^{\pm}(c) = \E\left[ e^{c\, j^{\pm}} \right], \qquad m^{\pm} = \ell^{\pm}(1) - 1,
\end{equation}
so that $\ell^{\pm}(c) = \mathcal{L}^{(\pm)}(-c)$ in the Laplace-transform notation of \cite{LiuPackhamSepp2026}. The intensities $\lambda^{\pm}_{t}$ do not clash with the market prices of risk $\lambda_{0}$, $\lambda_{1}$ of the volatility model, which do not appear in this section. The economy is that of Section \ref{sec:model}: the index in unit supply, the money market account at rate $r$, derivatives in zero net supply, and a representative CRRA agent with horizon $T$. The drift $\mu_{t}$ is endogenous. At the fraction $x_{t}$ in the index, wealth follows
\begin{equation}\label{eq:hawkes_wealth}
\frac{d\Pi_{t}}{\Pi_{t-}} = \left[ r + x_{t}(\mu_{t}-r) - x_{t}\lambda^{+}_{t}m^{+} - x_{t}\lambda^{-}_{t}m^{-} \right] dt + x_{t}\sigma\, dW_{t} + x_{t}(e^{j^{+}}-1)\, dN^{(1)}_{t} + x_{t}(e^{j^{-}}-1)\, dN^{(2)}_{t}.
\end{equation}

\subsection{HJB equation and equilibrium premia}

The value function $J(t,\Pi,\lambda^{+},\lambda^{-})$ satisfies the HJB equation
\begin{equation}\label{eq:hawkes_hjb}
\begin{split}
0 = \max_{x} \Big\{ &\partial_{t}J + \left[ r + x(\mu-r) - x\lambda^{+}m^{+} - x\lambda^{-}m^{-} \right]\Pi J_{\Pi} + \frac{1}{2}x^{2}\sigma^{2}\Pi^{2}J_{\Pi\Pi} \\
&+ \kappa^{+}(\theta^{+}-\lambda^{+})J_{\lambda^{+}} + \kappa^{-}(\theta^{-}-\lambda^{-})J_{\lambda^{-}} \\
&+ \lambda^{+} \E\left[ J\!\left(t, \left(1+x(e^{j^{+}}-1)\right)\Pi,\, \lambda^{+}+\beta_{11}j^{+},\, \lambda^{-}+\beta_{21}j^{+}\right) - J \right] \\
&+ \lambda^{-} \E\left[ J\!\left(t, \left(1+x(e^{j^{-}}-1)\right)\Pi,\, \lambda^{+}+\beta_{12}j^{-},\, \lambda^{-}+\beta_{22}j^{-}\right) - J \right] \Big\}.
\end{split}
\end{equation}
We apply the CRRA ansatz with an exponential-affine coefficient in the intensities:
\begin{equation}\label{eq:hawkes_ansatz}
J = g(t,\lambda^{+},\lambda^{-})\, \frac{\Pi^{\gamma}}{\gamma}, \qquad
g = \exp\left\{ q_{0}(t) + \psi^{+}(t)\,\lambda^{+} + \psi^{-}(t)\,\lambda^{-} \right\},
\end{equation}
where $\psi^{\pm} = \partial_{\lambda^{\pm}} \ln g$ are the risk-premium loadings on the two intensities, the exact analogues of $\psi = g_{\sigma}/g$ in the volatility model. The loadings define the effective jump tilts and the compensator gaps
\begin{equation}\label{eq:hawkes_chi}
\chi^{\pm} = (\gamma-1) + \psi^{+}\beta_{1i} + \psi^{-}\beta_{2i}, \qquad
\Delta^{\pm} = \ell^{\pm}(1+\chi^{\pm}) - \ell^{\pm}(\chi^{\pm}),
\end{equation}
with $i=1$ for positive and $i=2$ for negative jumps. Differentiating \eqref{eq:hawkes_hjb} in $x$ and evaluating at the market-clearing allocation $x^{*}=1$, where $\left(1+x(e^{j}-1)\right)^{\gamma-1}$ becomes $e^{(\gamma-1)j}$, yields the equilibrium equity premium:
\begin{equation}\label{eq:hawkes_premium}
\mu_{t} - r = (1-\gamma)\sigma^{2} + \lambda^{+}_{t}\left( m^{+} - \Delta^{+} \right) + \lambda^{-}_{t}\left( m^{-} - \Delta^{-} \right).
\end{equation}
In the jump-risk-premia notation of \cite{LiuPackhamSepp2026}, where the premia are the differences between risk-neutral and statistical jump compensators, Eq \eqref{eq:hawkes_premium} reads $\mu-r = (1-\gamma)\sigma^{2} - \gamma^{+}_{t} - \gamma^{-}_{t}$: the equity premium is the diffusive premium minus the two jump premia. The equilibrium thereby predicts that the statistical drift is affine in the intensities rather than constant.

Substituting \eqref{eq:hawkes_premium} back into \eqref{eq:hawkes_hjb} at $x^{*}=1$ reduces the HJB equation to a linear equation for $g$. The exponential-affine form \eqref{eq:hawkes_ansatz} solves this equation exactly, because the generator maps exponential-affine functions into themselves with affine coefficients. This contrasts with the volatility model, where a finite log-polynomial ansatz is a local closure. The loadings solve the system, in time to horizon $\tau = T-t$,
\begin{equation}\label{eq:hawkes_odes}
\begin{split}
\partial_{\tau}\psi^{+} &= -\kappa^{+}\psi^{+} + (1-\gamma)\,\ell^{+}(1+\chi^{+}) + \gamma\, \ell^{+}(\chi^{+}) - 1, \qquad \psi^{+}(0)=0,\\
\partial_{\tau}\psi^{-} &= -\kappa^{-}\psi^{-} + (1-\gamma)\,\ell^{-}(1+\chi^{-}) + \gamma\, \ell^{-}(\chi^{-}) - 1, \qquad \psi^{-}(0)=0,\\
\partial_{\tau}q_{0} &= \gamma r + \frac{1}{2}\gamma(1-\gamma)\sigma^{2} + \kappa^{+}\theta^{+}\psi^{+} + \kappa^{-}\theta^{-}\psi^{-}, \qquad q_{0}(0)=0.
\end{split}
\end{equation}
The system belongs to the same class as the transform ODEs of the moment generating function in \cite{LiuPackhamSepp2026}, a moment function of an affine argument plus a linear term, so the same Runge-Kutta machinery applies. The moment functions in \eqref{eq:hawkes_odes} are finite provided the tilted exponents stay inside the domains of $\ell^{\pm}$, which is the equilibrium analogue of the admissibility condition on $(\chi^{+},\chi^{-})$ in \cite{LiuPackhamSepp2026}.

\subsection{Identification of the measure change}

For the terminal-wealth value function, the pricing kernel is the marginal utility of wealth, $\zeta_{t} = J_{\Pi}(t,\Pi_{t},\lambda^{+}_{t},\lambda^{-}_{t}) = g\, \Pi_{t}^{\gamma-1}$. Combining the equations \eqref{eq:hawkes_premium} and \eqref{eq:hawkes_odes} shows that the drift of $\zeta_{t}$ equals $-r\zeta_{t}$, so the money market account is priced consistently and the density process
\begin{equation}\label{eq:hawkes_kernel}
M_{t} = e^{r t}\, \frac{\zeta_{t}}{\zeta_{0}}
\end{equation}
is a $\PP$-martingale defining $\QQ$. On a positive jump, $M_{t}$ multiplies by $e^{(\gamma-1)j^{+}}\, g(\lambda^{+}+\beta_{11}j^{+}, \lambda^{-}+\beta_{21}j^{+})/g = e^{\chi^{+} j^{+}}$, and on a negative jump by $e^{\chi^{-} j^{-}}$. The Brownian component carries the constant Girsanov shift $\varphi = (1-\gamma)\sigma$. The induced risk-neutral dynamics are therefore exactly the $\QQ$-family of \cite{LiuPackhamSepp2026}: the intensities scale as $\lambda^{\pm,\QQ}_{t} = \ell^{\pm}(\chi^{\pm})\,\lambda^{\pm}_{t}$ and the jump-size densities are exponentially tilted, $\varpi^{\pm,\QQ}(j) \propto e^{\chi^{\pm} j}\, \varpi^{\pm}(j)$. The identification with their density-process parametrization is
\begin{equation}\label{eq:hawkes_identification}
\xi^{+} = \xi^{-} = \gamma - 1, \qquad q_{1}^{\pm} = \psi^{\pm}, \qquad \varphi(t) = (1-\gamma)\sigma,
\end{equation}
so their internal intensity loadings correspond to the value-function loadings and their direct tilts correspond to the common marginal-utility exponent. Their state-dependent Brownian premium reduces to a constant under the equilibrium drift. The Brownian price of risk carries no hedging correction, in contrast with $\lambda_{0}$ in the volatility model, because the Brownian shock does not move the investment opportunity set.

Three consequences follow. First, a single risk-aversion parameter restricts the direct tilts to a common value, $\xi^{+}=\xi^{-}=\gamma-1$, while the calibration in \cite{LiuPackhamSepp2026} leaves $(\chi^{+},\chi^{-})$ free. Consistency of the calibrated tilts with a common $\xi$ is therefore a direct test of single-parameter CRRA pricing. Rejections measure tail-specific pricing beyond the preferences of the representative agent. Second, the effective tilts $\chi^{+}$ and $\chi^{-}$ differ from each other even under a single $\gamma$, through the loadings $\psi^{\pm}$ and the asymmetric excitation matrix in \eqref{eq:hawkes_chi}, so part of any measured asymmetry is preference-free clustering compensation. Third, the limits behave correctly: for $\gamma \to 1$ the loadings vanish and $\QQ = \PP$, and for log utility, $\gamma \to 0$, the loadings vanish and $\chi^{\pm} = -1$, the pure num\'{e}raire tilt with Esscher parameter one in the sense of \cite{GerberShiu1994}, with $\varphi = \sigma$.

\section{Equilibrium premia in the regime-switching log-normal SV model}\label{sec:regime}

We combine the quadratic-drift LogSV model with the transition-tied jump model of
\cite{Sepp2026gba}.  Regime 1 is growth and regime 2 is stress.  A transition
$1\to2$ is a crash and $2\to1$ is a recovery; there is no separate Poisson jump
process.  This distinction is important because the pricing kernel changes the
transition clock and its mark distribution jointly.

\subsection{Physical dynamics and return convention}

Let $N^{[ij]}(dt,dz)$ be the marked transition measure with $\PP$-compensator
$1_{\{\chi_{t-}=i\}}\lambda^{[ij]}F^{[i]}(dz)dt$, and put $h(z)=e^{z}-1$.
The compensated-return convention of \cite{Sepp2026gba} is
\begin{equation}\label{eq:regime_sv}
\begin{split}
\frac{dS_t}{S_{t-}}={}&\bar\mu^{[i]}dt+\sigma_t dW_t^{(0)}
 +\int h(z)\left(N^{[ij]}(dt,dz)-\lambda^{[ij]}F^{[i]}(dz)dt\right),\\
d\sigma_t={}&b^{[i]}(\sigma_t)dt+\beta^{[i]}\sigma_t dW_t^{(0)}
 +\varepsilon^{[i]}\sigma_t dW_t^{(1)},\\
b^{[i]}(\sigma)={}&(\kappa_1^{[i]}+\kappa_2^{[i]}\sigma)
(\theta^{[i]}-\sigma),
\end{split}
\end{equation}
The adjacent parenthesized factors in the last line are multiplied.
Volatility is continuous when the regime changes; the parameters, including its
attractor, switch after the price jump.  Equivalently, the raw between-transition
price drift is
\begin{equation}\label{eq:regime_drift_convention}
\mu^{[i]}=\bar\mu^{[i]}-\lambda^{[ij]}\alpha^{[i]},
\qquad \alpha^{[i]}=\E[h(J^{[i]})]=\ell^{[i]}(1)-1.
\end{equation}
The crash and recovery marks follow the paper's mean convention,
\begin{equation}\label{eq:regime_jump_laws}
J^{[1]}\sim-\operatorname{Exp}(1/\eta^{[1]}),\qquad
J^{[2]}\sim+\operatorname{Exp}(1/\eta^{[2]}),
\end{equation}
so $\eta^{[i]}$ is the mean absolute log jump, not the exponential rate, and
\begin{equation}\label{eq:regime_jump_mgf}
\ell^{[1]}(c)=\frac{1}{1+\eta^{[1]}c},\qquad
\ell^{[2]}(c)=\frac{1}{1-\eta^{[2]}c}.
\end{equation}

Let $g_i^H(t,\sigma)$ be the value coefficient for one fixed representative-agent
terminal date $H$.  With
$\psi_i^H=\partial_\sigma\log g_i^H$,
$\rho_i^H=g_j^H/g_i^H$, and
$\Delta_i=\ell^{[i]}(\gamma)-\ell^{[i]}(\gamma-1)$, the market-clearing
condition gives the two equivalent premium formulas
\begin{equation}\label{eq:regime_premium}
\begin{split}
\mu^{[i]}-r={}&(1-\gamma)\sigma^2-\beta^{[i]}\sigma^2\psi_i^H
-\lambda^{[ij]}\rho_i^H\Delta_i,\\
\bar\mu^{[i]}-r={}&(1-\gamma)\sigma^2-\beta^{[i]}\sigma^2\psi_i^H
+\lambda^{[ij]}\left(\alpha^{[i]}-\rho_i^H\Delta_i\right).
\end{split}
\end{equation}
The first line prices the raw between-transition drift; the second uses the total
expected-return convention reported in the allocation paper.  Mixing the two
conventions would omit the physical jump compensator.

\subsection{Coupled equilibrium coefficient and log-quadratic closure}

Substitution of \eqref{eq:regime_premium} into the HJB equation gives the linear
coupled PDE
\begin{equation}\label{eq:regime_system}
g_{i,t}^H+b^{[i]}g_{i,\sigma}^H
+\frac12\vartheta_i^2\sigma^2g_{i,\sigma\sigma}^H
+\left[\gamma r+\frac12\gamma(1-\gamma)\sigma^2\right]g_i^H
+\lambda^{[ij]}\left(\Lambda_i g_j^H-g_i^H\right)=0,
\end{equation}
where $g_i^H(H,\sigma)=1$,
$\vartheta_i^2=(\beta^{[i]})^2+(\varepsilon^{[i]})^2$, and
\begin{equation}\label{eq:regime_Lambda}
\Lambda_i=(1-\gamma)\ell^{[i]}(\gamma)
+\gamma\ell^{[i]}(\gamma-1).
\end{equation}
After removing the rate term and writing $h=H-t$, its Feynman--Kac representation is
\begin{equation}\label{eq:regime_FK}
\widehat g_i(h,\sigma)=\E_i^{\PP}\left[
\exp\left\{\frac12\gamma(1-\gamma)\int_0^h\sigma_s^2ds\right\}
\prod_{0<s\le h:\,\Delta\chi_s\ne0}\Lambda_{\chi_{s-}}
\,\middle|\,\sigma_0=\sigma\right].
\end{equation}
Every derivative maturity $T\le H$ uses this same $g^H$.  Re-solving the agent
problem with $H=T$ for each option would generate maturity-specific probability
measures rather than one equilibrium measure.

For $v_i=\sigma-\theta^{[i]}$ we use the published-FIRST, degree-two log expansion
\begin{equation}\label{eq:regime_ansatz}
\widehat g_i(h,\sigma)=\exp B_i(h,v_i),\qquad
B_i=a_i^{(0)}+a_i^{(1)}v_i+a_i^{(2)}v_i^2.
\end{equation}
At the same physical volatility, $v_j=v_i+\delta_i$ with
$\delta_i=\theta^{[i]}-\theta^{[j]}$.  Hence
$D_i^B=B_j(h,v_i+\delta_i)-B_i(h,v_i)=d_0^i+d_1^iv_i+d_2^iv_i^2$, where
\begin{equation}\label{eq:regime_dcoef}
\begin{split}
d_0^i&=a_j^{(0)}+a_j^{(1)}\delta_i+a_j^{(2)}\delta_i^2-a_i^{(0)},\\
d_1^i&=a_j^{(1)}+2a_j^{(2)}\delta_i-a_i^{(1)},\qquad
d_2^i=a_j^{(2)}-a_i^{(2)}.
\end{split}
\end{equation}
Taylor projection of $e^{D_i^B}$ through $v_i^2$ gives six nonlinear ODEs:
\begin{equation}\label{eq:regime_odes}
\begin{split}
\partial_h a_i^{(0)}={}&\frac12\vartheta_i^2(\theta^{[i]})^2s_i
+\frac12\gamma(1-\gamma)(\theta^{[i]})^2
+\lambda^{[ij]}\left(\Lambda_i e^{d_0^i}-1\right),\\
\partial_h a_i^{(1)}={}&-\bar\kappa_i a_i^{(1)}
+\vartheta_i^2\theta^{[i]}s_i
+2\vartheta_i^2(\theta^{[i]})^2a_i^{(1)}a_i^{(2)}
+\gamma(1-\gamma)\theta^{[i]}
+\lambda^{[ij]}\Lambda_i e^{d_0^i}d_1^i,\\
\partial_h a_i^{(2)}={}&-2\bar\kappa_i a_i^{(2)}-\kappa_2^{[i]}a_i^{(1)}
+\frac12\vartheta_i^2s_i
+4\vartheta_i^2\theta^{[i]}a_i^{(1)}a_i^{(2)}
+2\vartheta_i^2(\theta^{[i]})^2(a_i^{(2)})^2\\
&+\frac12\gamma(1-\gamma)
+\lambda^{[ij]}\Lambda_i e^{d_0^i}
\left(d_2^i+\frac12(d_1^i)^2\right),
\end{split}
\end{equation}
where $s_i=(a_i^{(1)})^2+2a_i^{(2)}$,
$\bar\kappa_i=\kappa_1^{[i]}+\kappa_2^{[i]}\theta^{[i]}$, and all coefficients
start from zero.  The PDE \eqref{eq:regime_system} is linear in $g_i$, but the
coefficient system \eqref{eq:regime_odes} is nonlinear.  The local single-regime
Riccati formula of Lemma \ref{lem:riccati} is the only literal closed form; the
regime solution is a semi-analytic finite-ODE approximation.

\subsection{Quadratic-preserving closure and relation to IJTAF}

There is a second, internally consistent approximation at the preceding Taylor
order.  It must be re-solved from the equilibrium PDE; it is not obtained by
deleting the cubic drift from the degree-two solution.  Define
\begin{equation}\label{eq:regime_linear_ansatz}
B_i^{\mathrm L}(h,v_i)=\widetilde a_i^{(0)}(h)
+\widetilde a_i^{(1)}(h)v_i,
\qquad
\widehat g_i^{\mathrm L}=\exp B_i^{\mathrm L}.
\end{equation}
At a common volatility, its regime difference is
\begin{equation}\label{eq:regime_linear_difference}
D_i^{\mathrm L}=B_j^{\mathrm L}(h,v_i+\delta_i)-B_i^{\mathrm L}(h,v_i)
=\widetilde d_0^i+\widetilde d_1^iv_i,
\end{equation}
where
$\widetilde d_0^i=\widetilde a_j^{(0)}+\widetilde a_j^{(1)}\delta_i
-\widetilde a_i^{(0)}$ and
$\widetilde d_1^i=\widetilde a_j^{(1)}-\widetilde a_i^{(1)}$.
Projection of \eqref{eq:regime_system} through degree one gives four coupled
nonlinear ODEs:
\begin{equation}\label{eq:regime_linear_odes}
\begin{split}
\partial_h\widetilde a_i^{(0)}={}&
\frac12\vartheta_i^2(\theta^{[i]})^2(\widetilde a_i^{(1)})^2
+\frac12\gamma(1-\gamma)(\theta^{[i]})^2
+\lambda^{[ij]}\left(\Lambda_i e^{\widetilde d_0^i}-1\right),\\
\partial_h\widetilde a_i^{(1)}={}&
-\bar\kappa_i\widetilde a_i^{(1)}
+\vartheta_i^2\theta^{[i]}(\widetilde a_i^{(1)})^2
+\gamma(1-\gamma)\theta^{[i]}
+\lambda^{[ij]}\Lambda_i e^{\widetilde d_0^i}\widetilde d_1^i,
\end{split}
\end{equation}
with zero initial conditions.  Unlike the scalar single-regime Riccati equation,
this regime-switching system is semi-analytic rather than a literal closed form.

\paragraph{CRRA consistency.}
The adjective ``log-linear'' describes the approximation to $\log \widehat g_i$;
it does not replace CRRA preferences by CARA preferences.  The utility function
$u(\Pi)=\Pi^\gamma/\gamma$, relative risk aversion $1-\gamma$, homothetic value
ansatz $J_i=g_i^H\Pi^\gamma/\gamma$, and market-clearing identities are unchanged.
Equations \eqref{eq:regime_linear_odes} solve exactly the degree-one projection of
the CRRA coefficient PDE, while substitution into the unprojected PDE leaves a
local residual beginning at $v_i^2$.  Closure $\mathrm L$ is therefore an
internally consistent \emph{approximate CRRA equilibrium}, not an exact global
solution of \eqref{eq:regime_system}.  Internal consistency additionally requires
the same $B_i^{\mathrm L}$ to determine $\psi_i^{\mathrm L}$, the Brownian market
prices of risk, $\rho_i^{\mathrm L}$, the marked transition compensator, the option
transform, and Monte Carlo dynamics.  Re-solving degree one and using it throughout
preserves that consistency; deleting only the cubic term from closure $\mathrm F$
does not.

The resulting equilibrium loading is state independent,
$\psi_i^{\mathrm L}=\widetilde a_i^{(1)}$, so both Brownian market prices of risk
are linear in volatility:
\begin{equation}\label{eq:regime_linear_mpr}
\lambda_{0,i}^{\mathrm L}
=\left[(1-\gamma)-\beta^{[i]}\widetilde a_i^{(1)}\right]\sigma,
\qquad
\lambda_{1,i}^{\mathrm L}
=-\varepsilon^{[i]}\widetilde a_i^{(1)}\sigma.
\end{equation}
Its induced volatility drift is exactly quadratic,
\begin{equation}\label{eq:regime_quadratic_Q_drift}
b_i^{\QQ,\mathrm L}
=-\bar\kappa_i v_i-\kappa_2^{[i]}v_i^2
-\beta^{[i]}(1-\gamma)(\theta^{[i]}+v_i)^2
+\vartheta_i^2\widetilde a_i^{(1)}(\theta^{[i]}+v_i)^2.
\end{equation}
The transition clock must simultaneously use
$\rho_i^{\mathrm L}=\exp(\widetilde d_0^i+\widetilde d_1^iv_i)$.
The direct Esscher tilt and the conditional risk-neutral jump-size law are the
same as under the log-quadratic closure; only the re-solved continuous loading and
regime-timing ratio differ.

There is no contradiction with IJTAF.  \cite{SeppRakhmonov2024} treat the linear
Brownian market prices of risk as model inputs, so their Girsanov drift adjustment
is quadratic.  Their degree-two
and degree-four coefficients belong to the option-pricing transform and never feed
back into the density kernel.  Here $B_i$ is an equilibrium value coefficient and
$\psi_i=\partial_\sigma B_i$ enters the Girsanov kernel itself.  A degree-two $B_i$
makes $\psi_i$ affine and $\sigma^2\psi_i$ cubic.  Degree one is precisely the
equilibrium approximation order at which the endogenous Brownian premia stay
linear and the quadratic-drift LogSV class remains invariant.

\begin{table}[htbp]
\centering
\small
\begin{tabular}{lll}
\toprule
Property & Log-linear closure $\mathrm L$ & Log-quadratic closure $\mathrm F$\\
\midrule
Degree of $B_i$ & 1 & 2\\
Equilibrium ODEs & 4 & 6\\
First omitted local power & $v_i^2$ & $v_i^3$\\
$\psi_i=\partial_\sigma B_i$ & constant & affine\\
Brownian premia & linear in $\sigma$ & quadratic in $\sigma$\\
$\QQ$ volatility drift & quadratic & cubic\\
$\log\rho_i$ & linear in $v_i$ & quadratic in $v_i$\\
Option-transform degree & 4 & 4\\
\bottomrule
\end{tabular}
\caption{The two internally consistent equilibrium closures compared in the
replication.  The label $\mathrm F$ denotes the full induced cubic dynamics, not
an additional Taylor truncation of the stochastic drift.}
\label{tab:regime_closures}
\end{table}

\subsection{Exact induced risk-neutral dynamics and the cubic term}

At a marked transition $i\to j$ the marginal-utility density jumps by
$\rho_i^H(t,\sigma)e^{(\gamma-1)z}$.  The full marked compensator under the
fixed-horizon equilibrium measure is therefore
\begin{equation}\label{eq:regime_Q_kernel}
\nu_{ij}^{\QQ,H}(t,\sigma,dz)=
\lambda^{[ij]}\rho_i^H(t,\sigma)e^{(\gamma-1)z}F^{[i]}(dz).
\end{equation}
It implies
\begin{equation}\label{eq:regime_Q}
q_{ij}^{\QQ,H}=\lambda^{[ij]}\rho_i^H\ell^{[i]}(\gamma-1),
\qquad
F_i^{\QQ}(dz)=\frac{e^{(\gamma-1)z}}{\ell^{[i]}(\gamma-1)}F^{[i]}(dz),
\end{equation}
and the exponential means become
\begin{equation}\label{eq:regime_Q_eta}
\eta_1^{\QQ}=\frac{\eta^{[1]}}{1-(1-\gamma)\eta^{[1]}},
\qquad
\eta_2^{\QQ}=\frac{\eta^{[2]}}{1+(1-\gamma)\eta^{[2]}}.
\end{equation}
Thus the chain is no longer independent of volatility under $\QQ$: its clock
depends on $t$ and $\sigma_t$ through $\rho_i^H$.

In compensated form, the exact induced dynamics are
\begin{equation}\label{eq:regime_Q_dynamics}
\begin{split}
\frac{dS_t}{S_{t-}}={}&r\,dt+\sigma_t dW_t^{(0),\QQ}
+\int h(z)\left(N^{[ij]}(dt,dz)-\nu_{ij}^{\QQ,H}(t,\sigma_t,dz)dt\right),\\
d\sigma_t={}&b_i^{\QQ,H}(t,\sigma_t)dt
+\beta^{[i]}\sigma_t dW_t^{(0),\QQ}
+\varepsilon^{[i]}\sigma_t dW_t^{(1),\QQ},\\
b_i^{\QQ,H}(t,\sigma)={}&b^{[i]}(\sigma)
-\beta^{[i]}(1-\gamma)\sigma^2+\vartheta_i^2\sigma^2\psi_i^H(t,\sigma).
\end{split}
\end{equation}
For the log-quadratic equilibrium approximation,
$\psi_i^H=a_i^{(1)}+2a_i^{(2)}v_i$, and hence
\begin{equation}\label{eq:regime_cubic_drift}
b_i^{\QQ,H}=-\bar\kappa_i v_i-\kappa_2^{[i]}v_i^2
-\beta^{[i]}(1-\gamma)(\theta^{[i]}+v_i)^2
+\vartheta_i^2(\theta^{[i]}+v_i)^2
\left(a_i^{(1)}+2a_i^{(2)}v_i\right).
\end{equation}
The coefficient of $v_i^3$ is exactly
$2\vartheta_i^2a_i^{(2)}$.  The cubic term is absent from the primitive
$\PP$-dynamics but present in the full induced $\QQ$-dynamics whenever
$a_i^{(2)}\ne0$.  Calling its unmatched action a PDE residual does not remove it
from the stochastic model.  Deleting only this term while retaining the degree-two
$B_i$, Brownian loadings, and timing ratio would mix two different density kernels
and is not an equilibrium approximation.  Our full-cubic Monte Carlo retains
\eqref{eq:regime_cubic_drift}.  The quadratic-preserving Monte Carlo instead uses
the separately solved \eqref{eq:regime_linear_odes},
\eqref{eq:regime_quadratic_Q_drift}, and $\rho_i^{\mathrm L}$ everywhere.
The difference between the two induced drifts is
\begin{equation}\label{eq:regime_drift_closure_difference}
b_i^{\QQ,\mathrm F}-b_i^{\QQ,\mathrm L}
=\vartheta_i^2\sigma^2\left[
(a_i^{(1)}-\widetilde a_i^{(1)})+2a_i^{(2)}v_i\right].
\end{equation}
It is a total closure difference: even at $v_i=0$ the two re-solved linear
coefficients need not coincide.

\subsection{Semi-analytic option transform}

Let $X_t=\log(S_t/\mathcal B_t)$ be the log price relative to the money-market
account, let an option expire at $T\le H$, and set $u=T-t$.  For the transform
$\E^{\QQ}[e^{-\Phi(X_T-X_t)}]$, write
\begin{equation}\label{eq:regime_option_ansatz}
G_i(u,x,v_i;\Phi)=e^{-\Phi x}\exp C_i(u,v_i;\Phi),
\qquad C_i(0,v_i;\Phi)=0.
\end{equation}
At backward option time $s\in[0,u]$, the remaining agent horizon is
$h_A=H-T+s$.  Put
$D_i^C=C_j(s,v_i+\delta_i;\Phi)-C_i(s,v_i;\Phi)$ and let $\mathcal P_m$
denote Taylor projection through degree $m$.  For closure
$c\in\{\mathrm L,\mathrm F\}$, use $B_i=B_i^c$,
$D_i^B=B_j^c-B_i^c$, and $b_i^{\QQ,H}=b_i^{\QQ,c}$ consistently.
Direct collection of the corresponding full generator gives
\begin{equation}\label{eq:regime_option_ode}
\begin{split}
\partial_s C_i=\mathcal P_m\Big[&
\left(b_i^{\QQ,H}(h_A,v_i)-\Phi\beta^{[i]}\sigma^2\right)C_i'
+\frac12\vartheta_i^2\sigma^2\left(C_i''+(C_i')^2\right)
+\frac12\Phi(\Phi+1)\sigma^2\\
&+\lambda^{[ij]}e^{D_i^B(h_A,v_i)}
\left\{\ell^{[i]}(\gamma-1-\Phi)e^{D_i^C}
+\Phi\Delta_i-\ell^{[i]}(\gamma-1)\right\}\Big],
\end{split}
\end{equation}
where $\sigma=\theta^{[i]}+v_i$.  We use $m=4$, the degree-four SECOND log
transform of \cite{SeppRakhmonov2024}, with either equilibrium closure.  The identities
$C_i(\Phi=0)=C_i(\Phi=-1)=0$ verify normalization and the discounted-stock
martingale.  Fourier inversion uses the line $\Phi=-1/2+iy$.  The implementation
collects \eqref{eq:regime_option_ode} directly as polynomials.  Under closure
$\mathrm F$, matched degree-three and degree-four contributions from the cubic
drift are retained; only higher powers enter the residual.  Under closure
$\mathrm L$, the exact quadratic $\QQ$ drift enters the same degree-four transform.

\subsection{Equity illustration, implied skews, and verification}

Table \ref{tab:regime_calibration} uses the transition intensities and jump sizes
from the equity illustration of \cite{Sepp2026gba}.  The reported $-25\%$ crash and
$+15\%$ recovery are expected arithmetic transition returns, implying
$\eta^{[1]}=1/3$ and $\eta^{[2]}=3/23$.  The remaining response coefficients come
from the maintained SPY LogSV illustration identified by the chapter's source
provenance, combined with the paper's growth/stress volatility levels.  They are
not distributed as a package calibration.  This hybrid is an illustration, not a calibration.
\begin{table}[htbp]
\centering
\small
\begin{tabular}{lrr}
\toprule
Parameter & Growth / crash & Stress / recovery\\
\midrule
$\theta^{[i]}$ & $15.0\%$ & $22.5\%$\\
$\lambda^{[ij]}$ & $0.10$ & $1.00$\\
$\E[e^{J^{[i]}}-1]$ & $-25.0\%$ & $+15.0\%$\\
$\eta^{[i]}$ & $1/3$ & $3/23$\\
$\kappa_1^{[i]}$ & $2.6949$ & $2.6949$\\
$\kappa_2^{[i]}$ & $10.0107$ & $10.0107$\\
$\beta^{[i]}$ & $-1.5082$ & $-1.5082$\\
$\varepsilon^{[i]}$ & $0.8503$ & $0.8503$\\
\bottomrule
\end{tabular}
\caption{Parameters used in the equity illustration.  The fixed agent horizon is
$H=3$ years, $\gamma=-0.5$ (relative risk aversion $1.5$), $r=0$, the initial
regime is growth, and $\sigma_0=15\%$.}
\label{tab:regime_calibration}
\end{table}
For the paper's crash law, the required jump moments exist for $\gamma>-2$, while
positivity of $\Lambda_1$ strengthens admissibility to $\gamma>-1$.  This is why
the old $\gamma=-2$ stress experiment cannot be combined with the equity jump
calibration.  At the baseline state the equilibrium tilt doubles the crash intensity
from $0.10$ to approximately $0.20$ and changes the expected crash from $-25\%$ to
$-40\%$; the recovery intensity changes from $1.00$ to approximately $0.84$ and its
mean jump from $+15\%$ to $+12.24\%$.

Figure \ref{fig:regime_smiles} uses the quadratic-preserving closure and shows the
resulting implied-volatility skews.  Higher risk aversion raises and steepens the
left wing because it makes crash marks more severe under $\QQ$.  The channel
experiment shows that the Esscher tail tilt is the dominant left-wing premium in
this calibration; diffusive and timing premia mainly shift the center and term
structure.  The package implements these channel scales as analytic,
non-equilibrium diagnostics.  The channel-isolation curves are martingale
counterfactuals, not separate equilibria, and non-unit scales do not define a
package Monte Carlo measure.  Starting in stress raises the right wing through the
recovery transition, while starting in growth produces the familiar index downside skew.
\begin{figure}[H]
\centering
\includegraphics[width=\textwidth]{regime_sv_smiles}
\caption{Quadratic-preserving equilibrium implied-volatility term structure,
risk-aversion sensitivity, premium-channel attribution, and initial-regime
dependence.  Strikes are normalized by the forward.}
\label{fig:regime_smiles}
\end{figure}

Figure \ref{fig:regime_premia} separates the mechanisms in the measure change.  The
timing loading is modest for the baseline, while the tail tilt moves crash severity
strongly with risk aversion.  Panel C compares the two separately solved $\QQ$
drifts with the physical drift; the quadratic curve is not obtained by deleting a
term from the full-cubic curve.
\begin{figure}[H]
\centering
\includegraphics[width=\textwidth]{regime_sv_premia}
\caption{Equilibrium transition, tail, volatility-drift, and value-loading channels.
Panels A and D use the quadratic-preserving closure; Panel C overlays both closures.}
\label{fig:regime_premia}
\end{figure}

Figure \ref{fig:regime_closure_comparison} makes the implied-volatility comparison
direct.  The three-month level curves are visually indistinguishable.  Across the
displayed strikes, the full-cubic minus quadratic-preserving difference is below
about $0.5$ volatility basis points at three months and $1.4$ volatility basis
points at one year for either initial regime.  Table
\ref{tab:regime_closure_skews} reports the fixed-moneyness skew
$\sigma_{\mathrm{IV}}(-0.20)-\sigma_{\mathrm{IV}}(+0.20)$; its closure difference
stays below one volatility basis point at all four maturities.  This is a numerical
finding for the displayed equity parameters, not a general identity.  Moreover,
the gap measures total closure sensitivity rather than the isolated cubic term,
because $\widetilde a_i^{(1)}$, $a_i^{(1)}$, and the transition timing ratios are
re-solved consistently.
\begin{figure}[H]
\centering
\includegraphics[width=\textwidth]{regime_sv_closure_comparison}
\caption{Comparison of consistent equilibrium closures.  Panels A and B show
three-month implied-volatility smiles under the log-linear equilibrium closure,
which preserves a quadratic risk-neutral volatility drift, and the log-quadratic
closure with its full induced cubic drift.  Panels C and D report
$10^4(\sigma_{\mathrm{IV}}^{\mathrm F}-\sigma_{\mathrm{IV}}^{\mathrm L})$ for four
maturities.  Physical parameters, preferences, jumps, agent horizon, and the
degree-four option transform are identical.}
\label{fig:regime_closure_comparison}
\end{figure}
\begin{table}[H]
\centering
\small
\input{\RegimeSvArtifactRoot/tables/regime_sv_closure_comparison_table.tex}
\caption{Fixed-moneyness implied-volatility skew comparison.  The skew columns are
percentage-volatility points; the last column is full cubic minus
quadratic-preserving in volatility basis points.}
\label{tab:regime_closure_skews}
\end{table}
The adopted replication separates production analytics from the independent
oracle.  The provisional package deep modules solve both equilibrium closures,
derive the induced $\QQ$ state, compute the option transforms, and run the
full-equilibrium $\QQ$ Monte Carlo for both initial regimes.  The frozen chapter
monolith and verifier supply the scalar Riccati residual, the exact
frozen-volatility matrix benchmark, and the independent physical-measure
Feynman--Kac Monte Carlo of \eqref{eq:regime_FK}.  The package checks the explicit
degree-one ODE identity and the exact quadratic and derived cubic drift
coefficients.  Following the event semantics of the GoalBasedAllocation
implementation, a regime jump changes the price and parameter state
simultaneously.  Under $\QQ$ the event probability is evaluated pathwise as
$1-e^{-q_{ij}^{\QQ,H}(t,\sigma)\,dt}$ because the equilibrium clock is state- and
time-dependent.  Package $\QQ$ Monte Carlo preserves the common Brownian increment
in price and volatility, uses the exact quadratic drift under closure $\mathrm L$,
and retains the exact induced cubic under closure $\mathrm F$.  For both, the
1,601-point Fourier grid is checked against a 3,201-point refinement.  The frozen
physical-measure comparison quantifies each polynomial-closure error rather than
mistaking it for Monte Carlo sampling error; package transform prices agree with
package $\QQ$ Monte Carlo within the stated confidence tolerance.
\begin{figure}[H]
\centering
\includegraphics[width=\textwidth]{regime_sv_validation}
\caption{Validation of both equilibrium closures and their degree-four option
transforms.  Panels A and B compare package $\QQ$ analytics with package $\QQ$
Monte Carlo; Panels C and D compare package equilibrium coefficients with the
frozen physical-measure Feynman--Kac oracle.  Error bars are 95\% Monte Carlo
intervals.}
\label{fig:regime_validation}
\end{figure}
\begin{table}[H]
\centering
\small
\input{\RegimeSvArtifactRoot/tables/regime_sv_validation_table.tex}
\caption{Selected report-quality package-$\QQ$ and frozen-physical-measure Monte
Carlo checks.}
\label{tab:regime_validation}
\end{table}
Production equilibrium, transform, and $\QQ$ Monte Carlo analytics are provided by
the provisional deep modules
\path{stochvolmodels.models.regime_logsv},
\path{stochvolmodels.models.regime_logsv_simulation}, and
\path{stochvolmodels.pricers.regime_switch_logsv_pricer}.  The frozen
\path{regime_switch_logsv.py} monolith and
\path{verify_regime_sv_equilibrium.py} remain independent oracles, not the
production pricing path.  Their immutable source hashes and the package adoption
commits are recorded in \path{source_provenance.json} and
\path{acceptance_manifest.json}.  From the repository root, the exact canonical
regeneration command is
{\scriptsize
\begin{verbatim}
python -m volatility_book.ch_lognormal_sv_risk_premia.generate_regime_sv_figures --profile canonical
\end{verbatim}
}
It writes the numerical payload, artifact manifest, figures, and tables beneath
the ignored canonical artifact root selected above.

\section{Discussion}\label{sec:discussion}

\subsection{Horizon dependence of the premia}
%% [NEW]
The loading depends on the time to the fixed representative-agent terminal date
$H$: it is $\widetilde a_i^{(1)}(H-t)$ under the log-linear closure and
$a_i^{(1)}(H-t)+2a_i^{(2)}(H-t)v_i$ under the log-quadratic closure.  Hence the
equilibrium premia are horizon-dependent and a priced derivative must expire no
later than $H$.  At $H$ the loading vanishes because $g^H(H,\sigma)=1$.  All
derivative maturities must nevertheless use the same $g^H$: changing $H$ with each
expiry changes the density process itself.  In the single-regime long-horizon limit
the log-linear loading converges to $a^{(1)}_{\infty}=r_{-}$ and the premium
coefficients become time-homogeneous.  The regime-switching implementation keeps
$H$ finite so this horizon convention remains explicit.

\subsection{Economic mechanism}
%% [AI-DRAFTED SKETCH - to be rewritten by the author per the project drafting rule:
%% the economic-mechanism passage is drafted by the author without AI assistance.]
For an equity index we have $\beta<0$ and $0<\gamma<1$, and Remark \ref{rem:branch} gives $a^{(1)}(\tau)>0$ for $\tau > 0$. Three consequences follow. First, the equity premium $\left[(1-\gamma) - \beta a^{(1)}\right]\sigma^{2}$ exceeds its constant-volatility value $(1-\gamma)\sigma^{2}$: the agent demands extra compensation for holding an asset whose volatility rises when its price falls, similar to being short a put option on the index. Second, $\bar\lambda_{1} = -\varepsilon a^{(1)}<0$, so volatility shocks orthogonal to returns carry a negative price of risk. Third, $\varphi<0$ implies that the volatility drifts higher under $\QQ$ than under $\PP$ through $\hat\kappa_{2} < \kappa_{2}$, which reproduces the negative variance risk premium documented by \cite{BakshiKapadia2003}.

\noindent\emph{Scope note.}  The preceding sign discussion assumes
$0<\gamma<1$.  It is not a description of the plotted regime illustration, whose
baseline is $\gamma=-0.5$ and whose displayed sensitivity values are non-positive.

\subsection{Variance premia versus skew premia}
%% [NEW] scope discussion: what diffusive premia can and cannot price.
An equivalent change of measure for a diffusion preserves the diffusion coefficients: the pricing kernel relocates drift and nothing else. The covariation $\sigma\beta(\sigma)$, which generates the skew, is therefore identical under $\PP$ and $\QQ$, and the short-dated at-the-money skew slope of the beta formulation, $\beta/2$, is the same for every choice of the risk premia. Preferences cannot make the market charge for skew in the volatility model: demand for downside protection has no parameter to flow into.

At finite maturity the premia do bend the skew, but only mechanically. A smaller $\hat\kappa_{2}$ makes the volatility more persistent under $\QQ$, so the skew decays more slowly with maturity. This effect vanishes as the maturity shrinks, cannot flip the sign anchored by $\beta$, and is driven by the same single combination $\beta\bar\lambda_{0}+\varepsilon\bar\lambda_{1}$ that sets the variance risk premium. In the volatility model, the finite-maturity skew premium is a shadow of the variance risk premium, not a separately priced object.

Jump models break exactly this constraint. The density process for a marked point process has functional freedom over the jump measure: the tilt in Section \ref{sec:hawkes} changes the frequency through $\ell^{\pm}(\chi^{\pm})$ and the severity through the tilted densities, per tail and independently across tails. The risk-neutral asymmetry then differs from the statistical asymmetry even at vanishing maturity, and it differs by an amount the market chooses. Demand for put protection maps onto $\chi^{-}$ and call overwriting supply maps onto $\chi^{+}$. \cite{BollerslevTodorov2011} document from the data side that the compensation for asymmetry sits in the jump tails.

We propose this division of labor as the organizing principle for risk-premia modeling. Diffusive volatility premia price the level and term structure of volatility, the variance risk premium, with one degree of freedom. Jump premia price asymmetry and tails, the skew and crash premia, with one degree of freedom per tail. The equilibrium spans both regimes with the single parameter $\gamma$, through $(\bar\lambda_{0}, \bar\lambda_{1})$ in the volatility model, through $\xi^{+}=\xi^{-}=\gamma-1$ in the Hawkes model, and through both channels at once in the regime-switching model of Section \ref{sec:regime}. Empirical rejection of the common tilt then measures the tail-specific pricing that the diffusive model has no room to express.

\subsection{Second-order accuracy}
%% [NEW]
The closed-form expansion \eqref{eq:affine} has log degree one and leaves powers beginning at $v^{2}$ uncontrolled.  In the terminology of \cite{SeppRakhmonov2024}, the FIRST affine expansion has log degree two and residual powers $v^{3}$ and $v^{4}$; the SECOND expansion has log degree four and residual powers $v^{5}$ through $v^{8}$.  The regime equilibrium system \eqref{eq:regime_odes} is the degree-two/FIRST closure, while the option transform \eqref{eq:regime_option_ode} uses degree four/SECOND.  Polynomial degree and Monte Carlo time-step convergence are checked separately in the replication.

\begin{thebibliography}{99}

\bibitem[Bakshi and Kapadia(2003)]{BakshiKapadia2003} Bakshi, G. and Kapadia, N. (2003). Delta-hedged gains and the negative market volatility risk premium. \textit{Review of Financial Studies}, 16(2), 527--566.

\bibitem[Bick(1990)]{Bick1990} Bick, A. (1990). On viable diffusion price processes of the market portfolio. \textit{Journal of Finance}, 45(2), 673--689.

\bibitem[Bollerslev and Todorov(2011)]{BollerslevTodorov2011} Bollerslev, T. and Todorov, V. (2011). Tails, fears, and risk premia. \textit{Journal of Finance}, 66(6), 2165--2211.

\bibitem[Cox et al.(1985)]{CIR1985} Cox, J. C., Ingersoll, J. E. and Ross, S. A. (1985). An intertemporal general equilibrium model of asset prices. \textit{Econometrica}, 53(2), 363--384.

\bibitem[Gerber and Shiu(1994)]{GerberShiu1994} Gerber, H. U. and Shiu, E. S. W. (1994). Option pricing by Esscher transforms. \textit{Transactions of the Society of Actuaries}, 46, 99--140.

\bibitem[He and Leland(1993)]{HeLeland1993} He, H. and Leland, H. (1993). On equilibrium asset price processes. \textit{Review of Financial Studies}, 6(3), 593--617.

\bibitem[Heston(1993)]{Heston1993} Heston, S. L. (1993). A closed-form solution for options with stochastic volatility with applications to bond and currency options. \textit{Review of Financial Studies}, 6(2), 327--343.

\bibitem[Karasinski and Sepp(2012)]{KarasinskiSepp2012} Karasinski, P. and Sepp, A. (2012). Beta stochastic volatility model. \textit{Risk}, October, 66--71.

\bibitem[Lewis(2000)]{Lewis2000} Lewis, A. L. (2000). \textit{Option Valuation under Stochastic Volatility}. Finance Press, Newport Beach.

\bibitem[Liu et al.(2026)]{LiuPackhamSepp2026} Liu, F., Packham, N. and Sepp, A. (2026). Jump risk premia in the presence of clustered jumps. Working paper.

\bibitem[Merton(1971)]{Merton1971} Merton, R. C. (1971). Optimum consumption and portfolio rules in a continuous-time model. \textit{Journal of Economic Theory}, 3(4), 373--413.

\bibitem[Sepp(2026)]{Sepp2026gba} Sepp, A. (2026). Dynamic mean--variance portfolio allocation under regime-switching jump-diffusions with absorbing barriers and distribution matching. Working paper.

\bibitem[Sepp and Rakhmonov(2024)]{SeppRakhmonov2024} Sepp, A. and Rakhmonov, P. (2024). Log-normal stochastic volatility model with quadratic drift. \textit{International Journal of Theoretical and Applied Finance}, 26(8), 2450003.

\end{thebibliography}

\end{document}
